\section{Two literal flat-source estimates and their direct repair}
\label{sec:general-flat-source-scope}

\paragraph{MA-GA-002 (counterexamples to two printed scalar inequalities).}
The literal scalar inequalities (5.3.3.21), printed/PDF
p.~256, and (5.3.3.23), p.~257, of Masullo's public draft
are false under their stated flat-source hypotheses.
This is a statement about those displayed estimates,
not a counterexample to the flat-space topological
identification in Lemma 5.3.3.3(3).  The latter estimates
needed in MA-GA-001 are proved directly by higher-order
Taylor bounds in (MA-GA.8).

\begin{proof}
In (5.3.3.21) take $\alpha=a=1$ and $\beta=0$.
The printed inequality specializes to
\[
 \sup_{x\ne0}(1+|\log|x||)|x|\,|f(x)|
       \leq \sup_x(1+|x|)|x|\,|f(x)|.
 \tag{MA-GA.33}
\]
Absolute values around the function do not weaken the
counterexample, since our functions are nonnegative.
Choose $\eta\in C_c^\infty((-1/4,1/4))$,
$0\leq\eta\leq1$, $\eta(0)=1$.  For $M\geq1$ put
\[
 f_M(x)=
 \begin{cases}
 x^{-1}\eta(\log x+M),&x>0,\\
 0,&x\leq0.
 \end{cases}
 \tag{MA-GA.34}
\]
Each $f_M$ is smooth, compactly supported away from zero,
and hence belongs to $\mathcal S_{00}(\mathbb R)$.
At $x=e^{-M}$ the left side of (MA-GA.33) is at least
$1+M$.  On the entire support, its right side is at most
$1+e^{-M+1/4}$.  Since $M>e^{-M+1/4}$ for $M\geq1$,
the inequality fails.  The ratio of the two sides is
unbounded as $M\to\infty$; thus merely inserting a
universal constant in that same one-seminorm inequality
would not repair it.  Using higher Schwartz derivatives,
as in (MA-GA.8), does repair the required continuity claim.
The failure does not depend on allowing derivative order
zero.  For $\alpha=a=\beta=1$ the same family satisfies
$(x\partial_x)(xf_M)=\eta'(\log x+M)$ and
$x^2f_M'=\eta'(\log x+M)-\eta(\log x+M)$.
At any fixed $t_0$ with $\eta'(t_0)\ne0$, the left
supremum is at least
$(1+M-|t_0|)|\eta'(t_0)|$, while the right one is at
most $(1+e^{-M+1/4})\|\eta'-\eta\|_\infty$.
This also fails for all sufficiently large $M$.

The first inequality in (5.3.3.23) asserts, for $|x|<A$,
\[
 \left|\frac{\varphi(x)}x\right|
       \leq\frac1{2A}\int_{-A}^{A}|\varphi'(y)|\,dy.
 \tag{MA-GA.35}
\]
Again the example is nonnegative at the evaluation points,
so it also disproves the version without the first absolute
value.  There is a single admissible flat function for which
this fails at arbitrarily small $A$.  Choose a nonnegative
smooth bump $\theta$, supported in $(-1/10,1/10)$,
increasing up to its value $\theta(0)=1$ and decreasing
afterward.  Its total variation is $2$.  Set
\[
 x_j=4^{-j},\qquad a_j=e^{-16^j},\qquad
 \varphi(x)=\sum_{j\geq1}a_j
                  \theta\!\left(\frac{x-x_j}{x_j}\right)
       \quad(x>0),\qquad \varphi(x)=0\quad(x\leq0).
 \tag{MA-GA.36}
\]
The supports $[0.9x_j,1.1x_j]$ are disjoint.  On the
$j$th support every derivative of order $k$ is bounded
by $C_k a_jx_j^{-k}$.  For every $N,k$ this is
$o(x_j^N)$ as $j\to\infty$.  The derivative functions
therefore extend continuously by zero at zero, with
all their difference quotients tending to zero.
Induction on $k$ proves that $\varphi$ is smooth and
flat at zero.  It is compactly supported and hence
belongs to $\mathcal S_{00}(\mathbb R)$ and, after
restriction, to the all-weight space
$\mathcal S_{(-\infty,\infty)}(\mathbb R^\times)$
used in that printed proof.

Take $A_j=2x_j$ and evaluate at $x=x_j$.
No support with index less than $j$ meets $[-A_j,A_j]$,
and every support with index at least $j$ is contained
there.  Consequently
\[
 \left|\frac{\varphi(x_j)}{x_j}\right|=\frac{a_j}{x_j},
 \qquad
 \frac1{2A_j}\int_{-A_j}^{A_j}|\varphi'(y)|\,dy
          =\frac1{2x_j}\sum_{m\geq j}a_m
          <\frac{a_j}{x_j}.
 \tag{MA-GA.37}
\]
For the strict inequality, $a_{m+1}/a_m\leq e^{-240}<1/2$
for $m\geq1$, so $\sum_{m>j}a_m<a_j$.
This disproves (MA-GA.35) along $A_j\downarrow0$ for
one fixed admissible function.

The correct estimate uses the segment from zero to the
actual point.  For $x>0$, flatness and the fundamental
theorem of calculus give
\[
 \left|\frac{\varphi(x)}x\right|
       \leq\frac1x\int_0^x|\varphi'(y)|\,dy
       \leq\frac{C_N}{N+1}x^N
       \quad\text{if }|\varphi'(y)|\leq C_Ny^N.
 \tag{MA-GA.38}
\]
The negative direction is identical after reflection.
Repeating this estimate for every derivative gives the
valid flat extension argument used in MA-GA-001.
\end{proof}

The two counterexamples are local source-audit corrections.
They do not alter the actual norm-cutoff theorem or supply
an additional arithmetic zero-exclusion condition.
